\chapter{威胁模型与系统总体设计}

本章首先界定系统的威胁模型与防护边界，明确"防什么、不防什么"；然后给出系统总体架构与数据流；接着建立掩模、噪声与视觉积分的核心数学模型；最后以"已实现、PoC 模拟、未来工作"三类对照表，清晰说明本原型对技术方案的覆盖范围。

\section{威胁模型与需求分析}

\subsection{攻击者能力假设}

本系统的防护对象是采用图像采集设备（以智能眼镜为代表）对显示器进行偷拍并自动识别的攻击者。结合背景技术中智能眼镜的实际工作方式，本文将攻击者能力划分为以下层次：

（1）被动抽帧攻击（Snapshots）：攻击者随机或定时截取单张高分辨率图片，送入 OCR 或目标检测识别。这是智能眼镜云端识别的典型方式。

（2）流式识别攻击（Streaming）：攻击者以视频流方式连续采集低分辨率帧，逐帧识别。这是实时翻译、场景理解类应用的典型方式。

（3）卷帘快门混合采样：手机等卷帘快门相机逐行曝光，单帧内可能混合多个子帧的片段。

（4）长曝光积分攻击：攻击者延长曝光时间，试图在单次曝光内积分多个子帧以还原原图。

（5）完整周期对齐平均攻击：攻击者主动连续录制、精确对齐帧相位，叠加覆盖一个或多个完整显示周期的所有子帧后做时域平均。这是能力最强的离线攻击。

\subsection{防护边界界定}

需要诚实指出的是，本系统作为基于线性时域叠加的方案，存在明确的能力边界。其数学本质是：人眼的"视觉无感"依赖完备性约束 $\sum_k I_k = I$ 的线性积分还原，而攻击者的时域平均使用的正是同一个求和运算。因此，凡是能被人眼积分还原的信息，理论上也能被覆盖完整周期的相机积分还原。这意味着第（5）类完整周期对齐平均攻击在纯线性方案下无法被严格防御。

基于上述分析，本系统的防护边界界定如下：

\textbf{有效防御}（本系统的目标威胁模型）：被动抽帧攻击、流式识别攻击、不足一个完整周期的多帧叠加、卷帘快门混合采样，以及配合反色帧/黑帧插入后的长曝光积分攻击。这些恰好覆盖了智能眼镜抽帧分析与流式识别两种真实采集策略。

\textbf{已知局限}（超出目标威胁模型）：攻击者主动连续录制、精确对齐相位、对覆盖完整周期的子帧做时域平均的离线攻击。本文不宣称防御该攻击，而是将其作为线性时域方案的原理性局限如实列出，并在第五章给出实验证据与缓解方向。

这一界定使本系统的安全声明严谨可辩护：在智能眼镜偷拍这一现实威胁场景下，攻击者通常处于"路过抓拍"或"实时流式"的被动采集状态，难以满足完整周期对齐平均所需的主动录制与精确同步条件。

\subsection{功能与非功能需求}

综合防护目标，本系统的需求归纳如表3-1所示。

\begin{table}[htbp]
  \centering
  \caption{系统需求分析}
  \begin{tabular}{cll}
    \toprule
    类别 & 需求项 & 说明 \\
    \midrule
    \multirow{4}{*}{功能需求}
      & 随机点阵掩模生成 & CSPRNG 驱动，满足完备性与互斥性 \\
      & 对抗噪声注入 & 时域互补，单帧增扰、积分抵消 \\
      & 子帧渲染合成 & 掩模、噪声、亮度补偿、反色帧 \\
      & 时序调度 & 伪随机置换、VBlank 同步、黑帧应急 \\
    \midrule
    \multirow{3}{*}{非功能需求}
      & 视觉无感 & FPI $<0.1$，$\Delta E<1.0$ \\
      & 防御有效 & 目标威胁下单帧 OCR 准确率接近 0 \\
      & 算法可验证 & 完备的单元测试与攻防实验 \\
    \bottomrule
  \end{tabular}
\end{table}

\section{系统总体架构}

本系统采用分层解耦的模块化架构，由四个核心功能模块和若干辅助模块组成。核心数据流方向为：原始帧缓冲 $\rightarrow$ 掩模生成器 $\rightarrow$ 噪声注入器 $\rightarrow$ 渲染驱动器 $\rightarrow$ 时序控制器 $\rightarrow$ 显示输出，如图3-1所示。

\begin{figure}[htbp]
  \centering
  \begin{tikzpicture}[
      node distance=5mm and 8mm,
      mod/.style={rectangle, draw, rounded corners=2pt, minimum height=11mm,
                  minimum width=26mm, align=center, font=\zihao{5}},
      arr/.style={-{Stealth[length=2.2mm]}, thick}]
    \node[mod] (src) {原始帧缓冲\\$I(x,y)$};
    \node[mod, right=of src] (mask) {掩模生成器\\ChaCha20};
    \node[mod, right=of mask] (noise) {噪声注入器\\FGSM/PGD};
    \node[mod, below=12mm of noise] (render) {渲染驱动器\\GPU/软件};
    \node[mod, left=of render] (timing) {时序控制器\\置换/VBlank};
    \node[mod, left=of timing] (disp) {显示输出\\$n$ 子帧序列};
    \draw[arr] (src) -- (mask);
    \draw[arr] (mask) -- (noise);
    \draw[arr] (noise) -- (render);
    \draw[arr] (render) -- (timing);
    \draw[arr] (timing) -- (disp);
    \node[draw, dashed, rounded corners, fit=(mask)(noise), inner sep=3mm,
          label={[font=\zihao{6}]above:子帧生成}] {};
  \end{tikzpicture}
  \caption{系统总体架构与数据流}
\end{figure}

各核心模块职责如下：掩模生成器负责由密钥和周期序号生成随机点阵掩模与置换序列；噪声注入器负责生成时域互补的对抗子噪声；渲染驱动器负责将掩模、噪声、亮度补偿在像素层面合成为子帧，支持 GPU 着色器渲染并在 GPU 不可用时自动回退到等价的软件渲染；时序控制器负责按伪随机置换顺序调度子帧输出、维护时序令牌、在渲染超时时插入黑帧。此外，系统还包含 HDR 亮度补偿、多显示器同步、视觉疲劳策略、配置持久化等辅助模块。

需要明确，本系统是\textbf{应用/算法层的概念验证原型}，子帧合成与时序调度在用户态完成，实时演示通过截取屏幕内容再在独立窗口中显示子帧的方式呈现效果；技术方案中描述的操作系统合成器层、GPU 驱动层透明注入（如 Windows Desktop Duplication、Linux KMS/DRM、macOS CoreDisplay）需要内核模块与代码签名，属于产品化阶段的未来工作，不在本原型实现范围内。

\section{核心数学模型}

\subsection{掩模的完备性与互斥性}

设原始图像为 $I(x,y)$，拆分因子为 $n$。掩模生成器为每个周期生成随机索引矩阵 $R(x,y) \in \{1,2,\dots,n\}$，由此构造 $n$ 个二值掩模：

\begin{equation}
M_k(x,y) =
\begin{cases}
1, & R(x,y) = k \\
0, & R(x,y) \ne k
\end{cases}
\end{equation}

第 $k$ 个子帧（暂不含噪声与补偿）为 $I_k(x,y) = I(x,y) \cdot M_k(x,y)$。掩模需满足两个约束：

\begin{equation}
\text{完备性：}\ \sum_{k=1}^{n} M_k(x,y) = 1; \qquad
\text{互斥性：}\ M_i(x,y) \cdot M_j(x,y) = 0\ (i \ne j)
\end{equation}

完备性保证每个像素在一个周期内恰好被点亮一次，从而 $\sum_k I_k = I$，人眼积分可无损还原；互斥性保证任一单帧只包含约 $1/n$ 的像素，信息高度碎片化。

\subsection{对抗噪声的时域互补}

噪声注入器在每个子帧上叠加对抗子噪声 $N_k(x,y)$，要求其满足时域互补约束：

\begin{equation}
\sum_{k=1}^{n} N_k(x,y) = 0
\end{equation}

例如 $n=4$ 时可取 $N_1=+\tfrac{1}{2}N_{\text{base}},\,N_2=-\tfrac{1}{2}N_{\text{base}},\,N_3=+\tfrac{1}{2}N_{\text{base}},\,N_4=-\tfrac{1}{2}N_{\text{base}}$。这样在单帧内噪声幅度为 $\tfrac{1}{2}\|N_{\text{base}}\|$，足以干扰 OCR；而 $n$ 帧积分后噪声完全抵消，不影响人眼观看。

\subsection{视觉积分还原}

综合掩模、噪声与亮度补偿，人眼在一个周期内积分得到的有效图像为：

\begin{equation}
E(x,y) = \frac{1}{n}\sum_{k=1}^{n}\big(I(x,y)\,M_k(x,y) + N_k(x,y)\big)\cdot B
= \frac{B}{n}\,I(x,y)
\end{equation}

其中 $B$ 为亮度补偿增益。由于 $\sum_k M_k = 1$ 且 $\sum_k N_k = 0$，积分结果正比于原图 $I(x,y)$，取背光增益 $B=n$ 即可恢复原始亮度。这一等式既是"人眼无感"的数学保证，也正是第五章所揭示的"完整周期可被还原"局限的根源所在——它对人眼和相机同时成立。

\section{技术方案实现范围对照}

为避免将本概念验证原型与技术方案中的产品级系统混淆，本文将技术要素划分为"已实现（真实算法）"、"PoC 模拟实现"和"未来工作"三类，如表3-2所示。后续章节及结论中对系统能力的所有声明均以此表为准。

\begin{table}[htbp]
  \centering
  \caption{技术方案实现范围对照}
  \begin{tabular}{cll}
    \toprule
    类别 & 技术要素 & 本原型中的状态 \\
    \midrule
    \multirow{5}{*}{已实现}
      & 随机点阵掩模 & ChaCha20+HKDF，完备/互斥/卡方检验 \\
      & 时域互补对抗噪声 & FGSM/PGD，$\sum N_k=0$ 精确成立 \\
      & 伪随机时序置换 & Fisher-Yates 无偏置换 \\
      & 反色帧/黑帧防御 & 长曝光、渲染超时应急 \\
      & 视觉无感与攻防评测 & FPI/$\Delta E$、多引擎 OCR/检测 \\
    \midrule
    \multirow{3}{*}{PoC 模拟}
      & 高刷新率时序 & 软件模拟 VBlank，外推高刷指标 \\
      & GPU 渲染链路 & moderngl 着色器，失败回退软件渲染 \\
      & 实时屏幕保护 & 截屏再显示，非驱动层透明注入 \\
    \midrule
    \multirow{2}{*}{未来工作}
      & 驱动/合成器层注入 & DXGI/KMS-DRM/CoreDisplay，需内核 \\
      & 真机高刷硬件验证 & 240/480Hz 实屏与受试者实验 \\
    \bottomrule
  \end{tabular}
\end{table}

\section{本章小结}

本章界定了系统的威胁模型，明确了"有效防御被动抽帧、流式识别、卷帘快门、插反色帧的长曝光，但不防完整周期对齐平均"的防护边界；给出了以四核心模块为中心的总体架构与数据流，并强调本系统是应用层概念验证原型；建立了掩模完备性/互斥性、噪声时域互补、视觉积分还原三个核心数学模型；最后以三类对照表清晰划分了技术要素的实现范围。这为后续模块实现与实验评测奠定了基础。
